% Options for packages loaded elsewhere
\PassOptionsToPackage{unicode}{hyperref}
\PassOptionsToPackage{hyphens}{url}
\PassOptionsToPackage{dvipsnames,svgnames,x11names}{xcolor}
%
\documentclass[
  letterpaper,
  DIV=11,
  numbers=noendperiod,
  oneside]{scrartcl}

\usepackage{amsmath,amssymb}
\usepackage{iftex}
\ifPDFTeX
  \usepackage[T1]{fontenc}
  \usepackage[utf8]{inputenc}
  \usepackage{textcomp} % provide euro and other symbols
\else % if luatex or xetex
  \usepackage{unicode-math}
  \defaultfontfeatures{Scale=MatchLowercase}
  \defaultfontfeatures[\rmfamily]{Ligatures=TeX,Scale=1}
\fi
\usepackage{lmodern}
\ifPDFTeX\else  
    % xetex/luatex font selection
\fi
% Use upquote if available, for straight quotes in verbatim environments
\IfFileExists{upquote.sty}{\usepackage{upquote}}{}
\IfFileExists{microtype.sty}{% use microtype if available
  \usepackage[]{microtype}
  \UseMicrotypeSet[protrusion]{basicmath} % disable protrusion for tt fonts
}{}
\makeatletter
\@ifundefined{KOMAClassName}{% if non-KOMA class
  \IfFileExists{parskip.sty}{%
    \usepackage{parskip}
  }{% else
    \setlength{\parindent}{0pt}
    \setlength{\parskip}{6pt plus 2pt minus 1pt}}
}{% if KOMA class
  \KOMAoptions{parskip=half}}
\makeatother
\usepackage{xcolor}
\usepackage[top=30mm,bottom=30mm,left=15mm,right=80mm,marginparwidth=60mm,marginparsep=10mm,heightrounded]{geometry}
\setlength{\emergencystretch}{3em} % prevent overfull lines
\setcounter{secnumdepth}{-\maxdimen} % remove section numbering
% Make \paragraph and \subparagraph free-standing
\makeatletter
\ifx\paragraph\undefined\else
  \let\oldparagraph\paragraph
  \renewcommand{\paragraph}{
    \@ifstar
      \xxxParagraphStar
      \xxxParagraphNoStar
  }
  \newcommand{\xxxParagraphStar}[1]{\oldparagraph*{#1}\mbox{}}
  \newcommand{\xxxParagraphNoStar}[1]{\oldparagraph{#1}\mbox{}}
\fi
\ifx\subparagraph\undefined\else
  \let\oldsubparagraph\subparagraph
  \renewcommand{\subparagraph}{
    \@ifstar
      \xxxSubParagraphStar
      \xxxSubParagraphNoStar
  }
  \newcommand{\xxxSubParagraphStar}[1]{\oldsubparagraph*{#1}\mbox{}}
  \newcommand{\xxxSubParagraphNoStar}[1]{\oldsubparagraph{#1}\mbox{}}
\fi
\makeatother


\providecommand{\tightlist}{%
  \setlength{\itemsep}{0pt}\setlength{\parskip}{0pt}}\usepackage{longtable,booktabs,array}
\usepackage{calc} % for calculating minipage widths
% Correct order of tables after \paragraph or \subparagraph
\usepackage{etoolbox}
\makeatletter
\patchcmd\longtable{\par}{\if@noskipsec\mbox{}\fi\par}{}{}
\makeatother
% Allow footnotes in longtable head/foot
\IfFileExists{footnotehyper.sty}{\usepackage{footnotehyper}}{\usepackage{footnote}}
\makesavenoteenv{longtable}
\usepackage{graphicx}
\makeatletter
\def\maxwidth{\ifdim\Gin@nat@width>\linewidth\linewidth\else\Gin@nat@width\fi}
\def\maxheight{\ifdim\Gin@nat@height>\textheight\textheight\else\Gin@nat@height\fi}
\makeatother
% Scale images if necessary, so that they will not overflow the page
% margins by default, and it is still possible to overwrite the defaults
% using explicit options in \includegraphics[width, height, ...]{}
\setkeys{Gin}{width=\maxwidth,height=\maxheight,keepaspectratio}
% Set default figure placement to htbp
\makeatletter
\def\fps@figure{htbp}
\makeatother

\usepackage{lastpage}
\usepackage{fancyhdr}
\usepackage{titling}
\pagestyle{fancy}
\fancyhead[L]{CMPE2150 PCB Project B}
\fancyhead[R]{Schematic Annotation}
\fancyfoot[C]{\thepage/\pageref*{LastPage}}
\setlength{\droptitle}{-2cm}
\KOMAoption{captions}{tablesignature}
\makeatletter
\@ifpackageloaded{tcolorbox}{}{\usepackage[skins,breakable]{tcolorbox}}
\@ifpackageloaded{fontawesome5}{}{\usepackage{fontawesome5}}
\definecolor{quarto-callout-color}{HTML}{909090}
\definecolor{quarto-callout-note-color}{HTML}{0758E5}
\definecolor{quarto-callout-important-color}{HTML}{CC1914}
\definecolor{quarto-callout-warning-color}{HTML}{EB9113}
\definecolor{quarto-callout-tip-color}{HTML}{00A047}
\definecolor{quarto-callout-caution-color}{HTML}{FC5300}
\definecolor{quarto-callout-color-frame}{HTML}{acacac}
\definecolor{quarto-callout-note-color-frame}{HTML}{4582ec}
\definecolor{quarto-callout-important-color-frame}{HTML}{d9534f}
\definecolor{quarto-callout-warning-color-frame}{HTML}{f0ad4e}
\definecolor{quarto-callout-tip-color-frame}{HTML}{02b875}
\definecolor{quarto-callout-caution-color-frame}{HTML}{fd7e14}
\makeatother
\makeatletter
\@ifpackageloaded{caption}{}{\usepackage{caption}}
\AtBeginDocument{%
\ifdefined\contentsname
  \renewcommand*\contentsname{Table of contents}
\else
  \newcommand\contentsname{Table of contents}
\fi
\ifdefined\listfigurename
  \renewcommand*\listfigurename{List of Figures}
\else
  \newcommand\listfigurename{List of Figures}
\fi
\ifdefined\listtablename
  \renewcommand*\listtablename{List of Tables}
\else
  \newcommand\listtablename{List of Tables}
\fi
\ifdefined\figurename
  \renewcommand*\figurename{Figure}
\else
  \newcommand\figurename{Figure}
\fi
\ifdefined\tablename
  \renewcommand*\tablename{Table}
\else
  \newcommand\tablename{Table}
\fi
}
\@ifpackageloaded{float}{}{\usepackage{float}}
\floatstyle{ruled}
\@ifundefined{c@chapter}{\newfloat{codelisting}{h}{lop}}{\newfloat{codelisting}{h}{lop}[chapter]}
\floatname{codelisting}{Listing}
\newcommand*\listoflistings{\listof{codelisting}{List of Listings}}
\makeatother
\makeatletter
\makeatother
\makeatletter
\@ifpackageloaded{caption}{}{\usepackage{caption}}
\@ifpackageloaded{subcaption}{}{\usepackage{subcaption}}
\makeatother
\makeatletter
\@ifpackageloaded{sidenotes}{}{\usepackage{sidenotes}}
\@ifpackageloaded{marginnote}{}{\usepackage{marginnote}}
\makeatother

\ifLuaTeX
  \usepackage{selnolig}  % disable illegal ligatures
\fi
\usepackage{bookmark}

\IfFileExists{xurl.sty}{\usepackage{xurl}}{} % add URL line breaks if available
\urlstyle{same} % disable monospaced font for URLs
\hypersetup{
  pdftitle={CMPE2150 PCB Project B},
  pdfauthor={AJ Armstrong},
  colorlinks=true,
  linkcolor={blue},
  filecolor={Maroon},
  citecolor={Blue},
  urlcolor={Blue},
  pdfcreator={LaTeX via pandoc}}


\title{CMPE2150 PCB Project B}
\usepackage{etoolbox}
\makeatletter
\providecommand{\subtitle}[1]{% add subtitle to \maketitle
  \apptocmd{\@title}{\par {\large #1 \par}}{}{}
}
\makeatother
\subtitle{Schematic Annotation}
\author{AJ Armstrong}
\date{09 Dec 2024}

\begin{document}
\maketitle

\renewcommand*\contentsname{Table of contents}
{
\hypersetup{linkcolor=}
\setcounter{tocdepth}{3}
\tableofcontents
}

\begin{tcolorbox}[enhanced jigsaw, opacityback=0, left=2mm, rightrule=.15mm, breakable, arc=.35mm, opacitybacktitle=0.6, colbacktitle=quarto-callout-tip-color!10!white, bottomrule=.15mm, toprule=.15mm, colback=white, colframe=quarto-callout-tip-color-frame, leftrule=.75mm, bottomtitle=1mm, toptitle=1mm, titlerule=0mm, title=\textcolor{quarto-callout-tip-color}{\faLightbulb}\hspace{0.5em}{Tip}, coltitle=black]

I'm going to stop reminding you to increment version numbers, check in
and push at major milestones, because you already do that reflexively.

\end{tcolorbox}

\subsubsection{Schematic Annotation}\label{schematic-annotation}

Now that we have the basic symbols and connections input, we're going to
modify this schematic, which currently represents the design of the
circuit in a generalized, abstract way to a schematic that also contains
(but does not always display) information about actual physical
components and characteristics

By the time we are finished, KiCAD will be satisfied that it knows
enough to help us lay out a PCB design, and our electrical rule check
(ERC) will be clean.

\subsubsection{Part 1 - Adding Values and
Footprints}\label{part-1---adding-values-and-footprints}

We will now visit each of the components in the schematic to correctly
set:

\begin{itemize}
\tightlist
\item
  Its \emph{value} (e.g.~\(R_1\) is \(2.2k\Omega\)), which you may have
  already selected when placing the symbol on the schematic.
\item
  Its \emph{footprint}, which is the actual, physical part description.
  It may depend on the manufacturer, type of part, or package.
\end{itemize}

In general, you can do both by double clicking on the symbol with the
selection tool or hover-selecting and hitting the \(E\) (Edit) key.
There, you can change the values in the appropriate fields. We will take
them in turn.

\begin{enumerate}
\def\labelenumi{\arabic{enumi}.}
\item
  We will first annotate each resistor (except the potentiometer).
  Confirm one has the correct value (for example, \(R_1\) should have a
  value of ``2.2k''). Choose an SMD-0805\sidenote{\footnotesize Surface Mount Device,
    0.08 inches by 0.05 inches.} footprint (search terms include ``R'',
  ``SMD'', and ``0805'') for that resistor. Continue and do the same for
  all ten standard resistors.

  \begin{tcolorbox}[enhanced jigsaw, opacityback=0, left=2mm, rightrule=.15mm, breakable, arc=.35mm, opacitybacktitle=0.6, colbacktitle=quarto-callout-tip-color!10!white, bottomrule=.15mm, toprule=.15mm, colback=white, colframe=quarto-callout-tip-color-frame, leftrule=.75mm, bottomtitle=1mm, toptitle=1mm, titlerule=0mm, title=\textcolor{quarto-callout-tip-color}{\faLightbulb}\hspace{0.5em}{Tip}, coltitle=black]

  Once you have done the first resistor, you can open the symbol table
  (\emph{Tools → Edit Symbol Fields}). There, as you can see, you can
  change fields for groups of components and manually edit values. When
  doing a lot of changes to values and footprints, this tool is
  invaluable. Note that, as you would expect, you can copy and paste
  values and footprints between symbols.

  \end{tcolorbox}
\item
  Leave the electrolytic caps (\(C_1\) and \(C_2\)) for later, but give
  \(C_3\) an SMD-0805 footprint as well. Note that the 0805 footprint
  for caps is not the same as that for resistors---in particular, it
  begins with a \emph{C} rather than an \emph{R}. Make sure you are in
  the correct section of the footprint selector for surface-mount
  capacitors.
\item
  Now do the diodes, excepting the LED. First, select one and edit
  (\(E\)) it. If the symbol isn't already \(1N4005\), you can select
  ``Change Symbol'' and choose a new library identifier, then click the
  library button to change it, and ``Change'' to enact the change. All
  should be \(1N4005G\) diodes, which you can update the ``Value'' field
  to display, with the same SMD-0805 footprint. Again, if you edit them
  in the \emph{Symbol Table} you can do all at once.

  \begin{tcolorbox}[enhanced jigsaw, opacityback=0, left=2mm, rightrule=.15mm, breakable, arc=.35mm, opacitybacktitle=0.6, colbacktitle=quarto-callout-note-color!10!white, bottomrule=.15mm, toprule=.15mm, colback=white, colframe=quarto-callout-note-color-frame, leftrule=.75mm, bottomtitle=1mm, toptitle=1mm, titlerule=0mm, title=\textcolor{quarto-callout-note-color}{\faInfo}\hspace{0.5em}{Note}, coltitle=black]

  Neatness always counts! As you change the display values of some of
  these symbols, take the time to move the labels on the schematic
  around to make it easy to read the values. In particular, avoid labels
  overlapping with symbols or wires. Remember you can rotate labels,
  too.

  \end{tcolorbox}
\item
  We'll use a standard through-hole bulb LED for the power indicator.
  Choose a footprint for a through-hole \(5.0mm\) LED. Once you have
  done so, amend the name (value) of the Footprint by appending
  ``\texttt{\_GREEN}'' to the name. Also add green to the value of the
  symbol itself so that it appears in the schematic.
\item
  Do the NPN transistors: value is ``2N3904'', change to the appropriate
  symbol if necessary, and select an SOT-23\sidenote{\footnotesize Small-Outline
    Transistor \#23} footprint.
\item
  Do the same for the 2N3906 PNP transistors.
\item
  The LM7805 should be in a T0220\sidenote{\footnotesize Transistor Outline \#220, a
    standard through-hole package for transistors very common for
    voltage regulators, which are also three-pin.} package footprint.
\item
  The LM339 should be in a SOIC-14\sidenote{\footnotesize Small Outline Integrated
    Circuit - 14 pins. A standard surface-mount package for DIP-style
    chips. Similar to the ICs that you use in your breadboard, but
    significantly smaller and with gull-wing leads suitable for surface
    mount soldering and pick-and-place machines.} package footprint. Of
  course, a real LM339 actually has four comparators, and we only need
  one. In the symbol properties, make sure that ``Unit A'' is selected,
  which means the first of the four comparator units. Try changing it to
  ``B'' or ``C'' and note what happens to the pin numbers on the
  schematic before changing back to ``A''.
\item
  You should have already selected a \texttt{Conn\_01x03\_Pin} header
  row for the three output pins. Now choose a footprint that has a
  \(2.54mm\) (\(0.1^"\)) pitch between them.
\item
  That should set us up with the parts that are already in the KiCAD
  database.
\end{enumerate}

\subsubsection{Part 2 - Selecting Parts From a
Vendor}\label{part-2---selecting-parts-from-a-vendor}

Some devices that we want to use aren't already in the KiCAD symbol
library, or we may not know what footprints are available. However,
device vendors are very good about supporting circuit and PCB designers,
and most major suppliers go out of their way to make sure it's easy to
use their parts in a new design.

\begin{enumerate}
\def\labelenumi{\arabic{enumi}.}
\item
  Let's start with the fuse. I'd rather not deal with mounting a fuse
  holder and replacing expendable fuses. Let's go to a \emph{polyfuse},
  which will heat up and automatically kill the power above a certain
  current, then gradually cool down and reset on its own. First, change
  the fuse's symbol to the generic polyfuse symbol in the symbol editor.
  Unfortunately, we don't have much in the way of real polyfuses in our
  default KiCAD library. Let's fix that.

  \begin{enumerate}
  \def\labelenumii{\alph{enumii}.}
  \tightlist
  \item
    Go to Digikey\sidenote{\footnotesize \url{https://www.digikey.ca/}} and search for
    the Littelfuse NANOSMDC016F-2, which is in a SMD-1206 form factor
    polyfuse\sidenote{\footnotesize Of course, in the real world for a real project,
      you would compare this part to other parts, check the availability
      and price, and probably order some.}.
  \item
    On that part's page, click on the link for the EDA/CAD Models. In
    the next page, click on ``Select Download Format''. Choose ``KiCAD
    (V6 and Later)'' and Download. Create a folder in your KiCad Project
    (I use one named \texttt{lib}) and copy the extracted files into it.
  \item
    Now you can import it into your KiCAD symbol library, possibly by
    just double clicking on it, depending on your platform and version.
    There should be a file in the downloaded symbol/model that describes
    how to import the \texttt{*.kicad\_sym} symbol using
    \emph{Preferences → Manage Symbol Libraries → Global Libraries} from
    the main (not Schematic Editor) window then browsing (little folder
    icon) and opening the file to add to your library.
  \item
    Perform a similar procedure to import the \texttt{*.kicad\_mod}
    model file Footprint using \emph{Preferences → Manage Footprint
    Libraries → Global Libraries}. Note that Footprint libraries are
    added by folder, so you won't be able to see the file---just select
    the folder that contains it.
  \item
    Now, go to your schematic and edit the polyfuse symbol, to select
    what you just imported. Try searching on the NANOSMDC\ldots{} part
    number for the symbol. The footprint you just imported can be found
    under the folder name that you selected (in my case, it show as a
    new dropdown named \texttt{lib)}.
  \item
    While you're in the symbol table, update the datasheet for the
    polyfuse to an appropriate datasheet URL on the web and change the
    value to \texttt{NANOSMDC016F-2}.
  \end{enumerate}
\item
  The two electrolytic capacitors are next. For \(C_1\) and \(C_2\),
  we'll select some through-hole barrel caps. Remember to set the value
  field for the symbols (syntax like ``220uF'' or ``220u'' will work).

  \begin{enumerate}
  \def\labelenumii{\alph{enumii}.}
  \tightlist
  \item
    Let's try Mouser\sidenote{\footnotesize \url{https://www.mouser.ca/}} this time.
    Search on ``electrolytic'', and select ``Aluminum Electrolytic
    capacitors - Radial Leaded''. This will give you a filter dialog.
    Select the capacitance value you need, ensure the voltage rating for
    the bigger cap is well above the peak voltage of its input (hint:
    what is the peak voltage of \(12V_{rms}\)?). Note that the filter
    has a greater than or equal to button so you can select above a
    certain voltage rather than exactly. The smaller cap must be able to
    manage more than \(5V\). From the available options, choose one that
    is in stock and reasonably priced, with the correct value. Ideally,
    you want to get both parts from the same manufacturer (you can
    filter by manufacturer, too).
  \item
    Whatever part you have chosen, go to its datasheet. First, confirm
    that it does come in the value and voltage we specified---you can't
    always trust the search filters. If you are satisfied, copy the URL
    of that datasheet to the capacitors datasheet field. Now, from the
    datasheet, find the actual size of the package. \emph{Make sure the
    size (in microfarads) of the cap you want is available in a
    particular footprint, as the datasheets will be for several
    different values.} In particular, we are interested in the barrel
    diameter and the pin pitch (distance between pins). In the footprint
    selector for the Cap, select \emph{Capacitor\_THT}\sidenote{\footnotesize Through-Hole
      Technology} and the the Radial Lead package that matches your
    selected device. Do this for both capacitors.
  \end{enumerate}
\item
  Our power switch will be a small slide switch. Go to one of our
  vendors and find a SPST through-hole slide switch---that can easily
  handle your 12 VAC suppy---that appeals to you. Adjust the symbol and
  footprint as needed. You may need to import the symbol and/or
  footprint from the vendor if you cannot find one in the existing
  library that matches. To find out, look at the data sheet for your
  selected switch, which will have detailed drawings. The most important
  values are the size and the spacing of the pins. You can look in the
  KiCAD library to find a matching footprint starting with
  \texttt{SW\_DIP\_SPSTx01\_Slide\_} for some generic footprints. You
  may also look at some of the standard slide switches and see if you
  can find it in the vendor library. If you can't match them, you can
  pick a slide switch whose footprint is available for download. (Hint:
  most footprints in the KiCAD library have links to the vendor.)
\item
  The potentiometer should be a through-hole vertical potentiometer.
  Select one from the \texttt{Potentiometer\_THT} footprints that has a
  knob that is adjustable without tools.
\item
  Finally, we need to do something about that \(12V_{RMS}\) supply.
  Obviously, we're not going to actually install an AC source onto our
  board. What is meant by that symbol is the wall-wart plug connected to
  our board to provide the AC power source. From the point of view of
  our board, that just appears as the barrel-jack connector our
  wall-wart will plug in to. So let's change it.

  \begin{enumerate}
  \def\labelenumii{\alph{enumii}.}
  \tightlist
  \item
    Edit the power supply symbol and select \emph{Change Symbol}. Choose
    a new library identifier to change it to a \texttt{Barrel\_Jack}.
    Make sure the pin (\(pin 1\)) is connected to the net you identified
    as \texttt{VAC} and the sleeve is connected to \texttt{GNDPWR}.
  \item
    Now select a footprint for the barrel jack: there is a DC barrel
    jack footprint.
  \end{enumerate}
\item
  Open the symbol table and confirm that all of the symbols have correct
  values and footprints. Visually clean up the schematic and labels.
  Increment the version number, check in and push.
\end{enumerate}

\subsubsection{Part 3 - Electrical Rules Check
(ERC)}\label{part-3---electrical-rules-check-erc}

At last, it is time to do the ERC. The notes below discuss how to
resolve the issues you are likely to encounter. If you have other
errors, try to resolve them. Just typing the error message and ``KiCAD''
into a search engine or AI is remarkably effective because of the large
KiCAD community online. Select \emph{Inspect → Electrical Rules
Checker}, and go to work.

\begin{enumerate}
\def\labelenumi{\arabic{enumi}.}
\tightlist
\item
  If you, as instructed, placed \texttt{PWR\_FLAG}s next to the
  \texttt{VCC} flag on the output of the Voltage regulator, move it to
  the input (pin 1) to get rid of the \emph{Input Power pin not driven
  by any Ouput Power pins} error. Although we are using it for
  \(V_{CC}\), pin 3 is already an output, so we do not need to use the
  flag to label it as getting an external output. Pin 1 is the actual
  input that needs the flag.
\item
  That no output for input power pin error is usually resolvable with
  power flags, as illustrated in the tutorial. Power flags are just a
  virtual output that labels a line as being an output to satisfy an
  input pin's need to be driven. You only need it when you have a device
  that has such a pin and the output to drive it is off-board.
\item
  If you get an error complaining about output power connected to an
  open emitter or open collector, you've likely got the wrong type of
  transistor---like a simulation transistor---make sure the symbol is a
  \texttt{2N3904} (or \texttt{-06}).
\item
  You will likely have an error that \(U_2\) has unplaced units. This is
  expected, as we only placed one of the four comparators on the
  \texttt{LM339}. In effect, we have three spare comparators that we
  aren't hooking up but we can't leave them out because they're part of
  the chip carrying the one they need. So, we'll place them.

  \begin{enumerate}
  \def\labelenumii{\alph{enumii}.}
  \tightlist
  \item
    Place three new comparators neatly in the white space below the
    circuit by copying the existing one and pasting them.

    \begin{enumerate}
    \def\labelenumiii{\alph{enumiii}.}
    \tightlist
    \item
      Edit them and change the \emph{Reference} to \texttt{U2} (the same
      as the first comparator) and the \(Unit\) to \texttt{B},
      \texttt{C}, etc. Note the different pin numbers.
    \item
      Place \emph{no-connection flags} (\$Q\$) on the inputs and outputs
      of the three comparators.
    \item
      Remember that we didn't have the power and ground connections in
      the schematic because they weren't really necessary (of course a
      chip is powered) and left out to reduce clutter? Well, the PCB
      Editor \emph{does} care how it's powered, so we will add them by
      placing \emph{Unit E} of the comparator. Copy a comparator one
      more time and change it to \texttt{U2E}. It'll turn into our
      missing power pins. Now, you can actually place those over the
      schematic symbol, but we don't need to. The PCB Editor just wants
      to know what to connect them to. Just place the two pins next to
      the spare comparators and place power symbols \texttt{VCC} and
      \texttt{GND} to connect them to.
    \end{enumerate}
  \end{enumerate}
\end{enumerate}

\subsubsection{Submission}\label{submission}

When you have completed all the above---including having no ERC errors
or warnings---and checked and reviewed your work, ask your instructor to
review your schematic with you. Make any changes they require. When you
are done submit the schematic.




\end{document}
